% Notes of the Montreal seminar, 1970 (cote 129, typed part): archive cote 129, site page 2 (PDF page 3).
% typescript. Draft compared: none.
% AI-transcribed by Claude Opus 5.5 (claude-opus-5-5, adaptive thinking, effort high), Message Batches API, batch mtl129_full (msgbatch_01AGnDKHwy6ATCM7ehUa6263), 2026-11-13; not yet checked by eye.
\scanpage{129/2}
1.  \uline{Rappels sur les vecteurs de Witt}.

Dans tout ce qui suit, $p$  désigne un nombre premier fixé
une fois pour toutes.

Le \uline{schéma des vecteurs de Witt} $\mathbb{W}$  et le \uline{schéma des vecteurs}
\uline{de Witt tronqués à l'ordre  $n$} $\mathbb{W}_n$  sont des schémas en anneaux, affines
sur $\mathbb{Z}$ qui sont définis comme suit:

Soit  $n \geq 1$  un entier et soit  $\mathbb{E}^n = \mathrm{Spec}\, \mathbb{Z}[T_1,T_2,\dots,T_n]$
d'espace affine de dimension  $n$  sur  $\mathbb{Z}$.  Le schéma  $\mathbb{E}^n$  représente
le foncteur  $A \mapsto A^n$  de la catégorie des schémas affines dans la
catégorie des ensembles; il est donc muni d'une structure naturelle de
schéma en anneaux.  On note  $\underline{O}^n$  le schéma  $\mathbb{E}^n$  muni de cette structure.
On démontre alors (c.f. J.-P. Serre, ``Corps locaux'' II §6) qu'il existe
une unique structure de schéma en anneaux sur  $\mathbb{E}^n$  telle que le morphisme
de schémas
\[
\phi = (\phi_1,\phi_2,\dots,\phi_n) \colon \mathbb{E}^n \to \underline{O}^n
\]
défini par
\[
\phi_i(x_1,x_2,\dots,x_n) = x_1^{p^{i-1}} + p\, x_2^{p^{i-2}} + p^2\, x_3^{p^{i-3}} + \dots + p^{i-1}\, x_i
\]
soit un homomorphisme d'anneaux.

Par définition, $\mathbb{W}_n$  est le schéma  $\mathbb{E}_n$  muni de cette structure
de schéma en anneaux.

\mspage{\hand{2}}
